\section*{Chapter \ref{ch:g10:national-accounts} --- National Accounts}
% figdata: figdata/grade-10/11-national-accounts/out/ (shares, gni, brenland, bren_2025); Brenland numbers invented.

\begin{solution}{exo:g10:national-accounts:1}
Consumption is largest in Kenya (77.0\,\%), \omterm{def:g10:national-accounts:investment}{investment} in China (41.1\,\%). \omterm{def:g10:national-accounts:net-exports}{Net exports} are negative in the United States ($-2.8$), India ($-1.6$) and Kenya ($-7.5$).
\end{solution}

\begin{solution}{exo:g10:national-accounts:2}
\omterm{def:g10:national-accounts:investment}{Investment}: $41.1 - 21.6 = 19.5$ points more in China. Consumption: $67.7 - 39.6 = 28.1$ points less.
\end{solution}

\begin{solution}{exo:g10:national-accounts:3}
$Y = 600 + 200 + 250 + 300 - 350 = 1000$; \omterm{def:g10:national-accounts:net-exports}{net exports} $300 - 350 = -50$.
\end{solution}

\begin{solution}{exo:g10:national-accounts:4}
Forest 40, sawmill $70 - 40 = 30$, furniture maker $150 - 70 = 80$. GDP $= 40 + 30 + 80 = 150$, the value of the tables bought by households.
\end{solution}

\begin{solution}{exo:g10:national-accounts:5}
$800 + 30 - 70 = 760$.
\end{solution}

\begin{solution}{exo:g10:national-accounts:6}
Year 1: $100 \times 500/500 = 100$. Year 2: $100 \times 540/515 \approx 104.85$. The deflator rose by about 4.85\,\%.
\end{solution}

\begin{solution}{exo:g10:national-accounts:7}
$100 \times 152000/197403 \approx 77.0$. Then $134.5/77.0 \approx 1.747$: prices of Brenland's output rose by about 74.7\,\% over the 25 years.
\end{solution}

\begin{solution}{exo:g10:national-accounts:8}
Exactly: $1.07/1.04 \approx 1.0288$, real growth about 2.88\,\%. Approximately: $7 - 4 = 3\,\%$.
\end{solution}

\begin{solution}{exo:g10:national-accounts:9}
$1000 - 150 = 850$. If more capital wears out each year (more of it, or of a kind that ages fast, like computers), depreciation can grow as fast as GDP or faster, and the
net product then stands still while GDP grows.
\end{solution}

\begin{solution}{exo:g10:national-accounts:10}
Before, households' consumption included 1200 of food; now it includes 3000 of meals, and the 1200 of food is an input of the shop. GDP rises by $3000 - 1200 = 1800$, the
shop's value added. The family eats the same food; what changed is that the cooking is now done for pay, and counted.
\end{solution}

\begin{solution}{exo:g10:national-accounts:11}
(a) $C$ rises and $M$ rises by the same amount: GDP unchanged. (b) $I$ rises: GDP rises. (c) A transfer, not a purchase: GDP unchanged. (d) A second-hand house is not new
production: GDP unchanged (except the agent's fee, a service). (e) The change in stocks counts as \omterm{def:g10:national-accounts:investment}{investment}: GDP rises by the value of the shoes.
\end{solution}

\begin{solution}{exo:g10:national-accounts:12}
The proposition is about the true totals. The published shares come from different sources (household surveys, \omterm{def:g10:national-accounts:investment}{investment} surveys, customs data), each with its own gaps
and errors, so the expenditure side need not add exactly to the GDP that the production side gives. The gap, about 3 points of GDP for Vietnam and 2 for Kenya, is a
statistical discrepancy. It means that each share is uncertain by an amount of that order, so small differences between countries should not be read too closely.
\end{solution}

\begin{solution}{exo:g10:national-accounts:13}
$103/100$: 3\,\%; $105/103$: about 1.94\,\%; $110/105$: about 4.76\,\%; $112/110$: about 1.82\,\%. The product $1.03 \times 1.0194 \times 1.0476 \times 1.0182 \approx 1.12$ equals
$112/100$, since each factor's denominator cancels the next one's numerator.
\end{solution}

\begin{solution}{exo:g10:national-accounts:14}
By expenditure, $C = 400$ and $G = 150 + 30 = 180$ (the schooling valued at its cost): GDP $= 580$. By production: bread 400, the firms making the materials 30, the school's own
value added 150: again 580. After the pay rise, the school is valued at $165 + 30 = 195$ and GDP becomes 595, although nothing more is taught.
\end{solution}

\begin{solution}{exo:g10:national-accounts:15}
Expenditure: $X = 100$ and $M = 70$, so the machine adds $100 - 70 = 30$. Production: the firm's value added is sales minus purchased inputs, $100 - 70 = 30$. The same 30.
\end{solution}

\begin{solution}{pb:g10:national-accounts:1}
Money in millions of dalens unless stated.

\textbf{1.} $225014 + 88484 + 78653 + 129777 - 128664 = 393264$.

\textbf{2.} $C$ 57.2\,\%, $I$ 22.5\,\%, $G$ 20.0\,\%, $X$ 33.0\,\%, $M$ 32.7\,\%.

\textbf{3.} 2024: $125104 - 126420 = -1316$. 2025: $129777 - 128664 = 1113$.

\textbf{4.} $C$ 2.48\,\%, $I$ 3.61\,\%, $G$ 3.61\,\%, $X$ 3.74\,\%, $M$ 1.78\,\%. Exports grew fastest.

\textbf{5.} $C$ $+5449$, $I$ $+3083$, $G$ $+2741$, $X$ $+4673$, $M$ $-2244$: in all $+13702$, the rise of GDP from \num{379562} to \num{393264}.

\textbf{6.} $393264/379562 \approx 1.0361$: 3.61\,\%.

\textbf{7.} $292390/286895 \approx 1.0192$: 1.92\,\%.

\textbf{8.} 2024: $100 \times 379562/286895 \approx 132.3$. 2025: $100 \times 393264/292390 \approx 134.5$.

\textbf{9.} $134.5/132.3 \approx 1.0166$: 1.66\,\%. Check: $1.0192 \times 1.0166 \approx 1.0361$.

\textbf{10.} $1.66/3.61 \approx 46\,\%$ of the nominal growth.

\textbf{11.} $9291/9235 \approx 1.0061$: 0.61\,\%.

\textbf{12.} 2024: $286895/9235 \approx 31.066$ thousand dalens, \dal{31066}. 2025: $292390/9291 \approx 31.470$, \dal{31470}.

\textbf{13.} $31470/31066 \approx 1.0130$: 1.30\,\%. Approximately $1.92 - 0.61 = 1.31\,\%$.

\textbf{14.} $393264/9291 \approx 42.327$ thousand: \dal{42327}.

\textbf{15.} $292390/248438 \approx 1.177$: 17.7\,\% in ten years, an average of $1.177^{1/10} \approx 1.0164$, 1.64\,\% a year.

\textbf{16.} $393264 + 4000 - 9000 = 388264$.

\textbf{17.} $393264 \times 0.88 \approx 346072$.

\textbf{18.} By a quarter of GDP, about \num{98316} million dalens.

\textbf{19.} Only if unpaid work grew at a different rate from measured GDP. One would need its value in both years: if, for example, more people bought services they used to
do at home, measured GDP grew faster than the total.

\textbf{20.} In 2025 Brenland's real GDP per person grew by 1.30\,\%. Of the 3.61\,\% rise in \omterm{def:g10:national-accounts:nominal-gdp}{nominal GDP}, 1.66 points, about 46\,\%, came from higher prices.
\end{solution}
